% This is part of Mes notes de mathématique
% Copyright (c) 2011-2020
%   Laurent Claessens
% See the file fdl-1.3.txt for copying conditions

%+++++++++++++++++++++++++++++++++++++++++++++++++++++++++++++++++++++++++++++++++++++++++++++++++++++++++++++++++++++++++++ 
\section{Groupes}
%+++++++++++++++++++++++++++++++++++++++++++++++++++++++++++++++++++++++++++++++++++++++++++++++++++++++++++++++++++++++++++

%---------------------------------------------------------------------------------------------------------------------------
\subsection{Définition, unicité du neutre}
%---------------------------------------------------------------------------------------------------------------------------

\begin{definition}[Groupe]      \label{DEFooBMUZooLAfbeM}
    Un \defe{groupe}{groupe} est un ensemble \( G\) muni d'une opération interne \( \cdot\colon G\times G\to G\) telle que
    \begin{enumerate}
        \item
            pour tous \( g\), \( h\), \( k\in G\), \( g\cdot(h\cdot k)=(g\cdot h)\cdot k\),
        \item
            il existe un élément \( e\in G\) tel que \( e\cdot g=g\cdot e=g\) pour tout \( g\in G\),
        \item
            pour tout \( g\in G\), il existe un élément \( h\in  G\) tel que \(g\cdot h=h\cdot g=e \).
    \end{enumerate}
\end{definition}

    Notons que nous avons écrit \( g\cdot h\) et non \( \cdot(g,h)\) comme une notation purement fonctionnelle nous l'aurait suggéré. Dans les exemples concrets, selon les cas, le loi de groupe appliquée à \( g\) et \( h\) sera notée tantôt \( g+h\), tantôt \( g\cdot h\) ou, le plus souvent pour un groupe générique, simplement \( gh\).

\begin{lemmaDef}[Unicités]  \label{LEMooECDMooCkWxXf}
    L'inverse et le neutre sont uniques, c'est-à-dire :
    \begin{enumerate}
        \item
            il existe un unique élément \( e\in G\) tel que \( e g=g e=g\) pour tout \( g\in G\),
        \item       \label{ITEMooOIWTooYqmMPP}
            pour tout \( g\in G\), il existe un unique élément \( h\in  G\) tel que \(g h=h g=e \).
    \end{enumerate}
    Le \( e\) ainsi défini est nommé \defe{neutre}{neutre!dans un groupe} de \( G\). Le \( h\) tel que \( g h=h g=e\) est nommé l'\defe{inverse}{inverse!dans un groupe} de \( g\) et est noté \( g^{-1}\).
\end{lemmaDef}

\begin{proof}
    Chaque point séparément.
    \begin{enumerate}
        \item
            Supposons que \( e_1\) et \( e_2\) vérifient la propriété. Nous avons pour tout \( g\in G\) : \( e_1g=ge_1=g\). En particulier pour \( g=e_2\) nous écrivons \( e_1e_2=e_2e_1=e_2\). Mais en partant dans l'autre sens : \( e_2g=ge_2=g\) avec \( g=e_1\) nous avons \( e_2e_1=e_1e_2=e_1\). En égalant ces deux valeurs de \( e_2e_1\) nous avons \( e_1=e_2\).

            Pour la suite de la preuve nous écrivons \( e\) l'unique neutre de \( G\).

        \item
            Supposons que \( k_1\) et \( k_2\) soient deux inverses de \( g\). On considère alors le produit \( k_1 g k_2 \). Puisque \(k_1 g = e \), on a \( k_1 g k_2 = e k_2 = k_2 \); mais, comme \(g k_2 = e \), on a aussi \( k_1 g k_2 = k_1 e = k_1 \). Le produit est donc à la fois égal à \( k_1 \) et à \( k_2 \), et donc \( k_1 = k_2 \).
    \end{enumerate}
\end{proof}

\begin{definition}\label{DEFGroupeAbelien}
    Un groupe \( G\) est \defe{abélien}{abélien!groupe} ou \defe{commutatif}{commutatif!groupe} si pour tous \( g\) et \( h\) dans \( G\), \( gh=hg\).
\end{definition}

%-----------------------------------
\subsection{Puissances dans un groupe}
%---------------------------------------

Voici une première définition de la fonction puissance. Il y en aura d'autres, de plus en plus générales. Voir le thème \ref{THEMEooBSBLooWcaQnR}.
\begin{definition}\label{DEFooGVSFooFVLtNo}
    Si \( G\) est un groupe (multiplicatif), si \( a\in G\) et si \( n\in \eN\), nous définissons \( a^n\) par récurrence :
    \begin{enumerate}
        \item
            \( a^0=1\) (le neutre),
        \item       \label{ITEMooOUIPooGjAgQb}
            \( a^{k+1}=a\cdot a^{k}\).
    \end{enumerate}
\end{definition}

Le lemme suivant dit que le point \ref{ITEMooOUIPooGjAgQb} de la définition \ref{DEFooGVSFooFVLtNo} aurait pu être écrit \( a^k\cdot a\) au lieu de \( a\cdot a^k\).
\begin{lemma}[\cite{MonCerveau}]        \label{LEMooWPARooYLZlzr}
    Si \( G\) est un groupe, si \( a\in G \) et si \( n\in \eN\) est non-nul, alors
    \begin{equation}
        a^n=a\cdot a^{n-1}=a^{n-1}\cdot a.
    \end{equation}
\end{lemma}
\begin{proof}
Pour \( n = 1 \), c'est trivial; pour \( n = 2 \), c'est facile. Pour des \( n \) plus grands, par récurrence, on suppose que \( a^{n-1} = a^{n-2}\cdot a \), et il vient 
\begin{align}
 a^n = a\cdot a^{n-1} &= a\cdot  a^{n-2}\cdot a  &\text{par hypothèse}\\
 & = a^{n-1}\cdot a &\text{par définition}.
 \end{align}
\end{proof}

%---------------------------------------------------------------------------------------------------------------------------
\subsection{Morphismes de groupes}
%---------------------------------------------------------------------------------------------------------------------------
\begin{lemma}       \label{LEMooBIBFooBHxFYC}
    Si \( G\) est un groupe et si \( h\in G\), alors les applications
    \begin{equation}
        \begin{aligned}
            L_h\colon G&\to G \\
            g&\mapsto hg 
        \end{aligned}
    \end{equation}
    et
    \begin{equation}
        \begin{aligned}
            R_h\colon G&\to G \\
            g&\mapsto gh 
        \end{aligned}
    \end{equation}
    sont des bijections (mais ne sont pas des morphismes\footnote{La définition suit ce lemme!} quand \(h \neq e \)).
\end{lemma}

\begin{proof}
    Ce ne sont, en effet, pas des morphismes, puisque \(e \) n'est pas envoyé sur \(e \), sauf cas bien particulier\dots\ Montrons à présent que ce sont tout de même des bijections.

    D'abord si \( L_h(g_1)=L_h(g_2)\), alors \( hg_1=hg_2\) et en multipliant à gauche par \( h^{-1}\) nous avons \( g_1=g_2\); donc \( L_h\) est injective. Ensuite \( L_h\) est surjective parce que si \( g\in G\), alors \( g=L_h(h^{-1} g)\).

    Pour l'application \( R_h\), la preuve est une simple adaptation.
\end{proof}

\begin{definition}
    Soit \(G,\ H\) deux groupes, et \(f \) une application de \(G \) dans \(H \). Si \(f \) respecte les structures de groupe, alors \(f \) est appelé un \defe{morphisme de groupes}{morphisme!de groupes} (ou un \emph{homomorphisme de groupes}). Respecter les structures de groupe signifie que:
    \begin{enumerate}
    \item \( f(e_G) = e_H \): le neutre de \( G \) est envoyé sur le neutre de \( H \);
    \item pour tous \(g, g' \in G,\ f(gg') = f(g)f(g')\): l'image d'un produit (dans \(G \)) est le produit (dans \( H \)) des images;
    \item pour tout \(g \in G,\ f(g^{-1}) = \left(f(g)\right)^{-1}\): l'image de l'inverse d'un élément est l'inverse de l'image de cet élément.
    \end{enumerate}
\end{definition}

\begin{definition}\label{DEFGroupesTypesdeMorphismes}
Un morphisme injectif s'appelle un \emph{monomorphisme}.

Un morphisme surjectif s'appelle un \emph{épimorphisme}.

Un morphisme bijectif s'appelle un \emph{isomorphisme}.

Les morphismes d'un groupe dans lui-même sont appelés des \emph{endomorphismes}. Quand ils sont bijectifs, on dit que ce sont des \emph{automorphismes}.
\end{definition}

Il y a plein d'autres choses à dire sur les groupes. On y reviendra entre autres au chapitre~\ref{CHAPGroupes}.

%+++++++++++++++++++++++++++++++++++++++++++++++++++++++++++++++++++++++++++++++++++++++++++++++++++++++++++++++++++++++++++ 
\section{Anneaux}
%++++++++++++++++++++++++++++++++++++++++++++++++++++++++++++++++++++++++++++++++++++++++++++++++++++++++++++++++++++++++++

\subsection{Définitions et les bases}
%--------------------------------------

\begin{definition}[Anneau\cite{Tauvel}]     \label{DefHXJUooKoovob}
    Un \defe{anneau}{anneau}\footnote{Nous faisons le choix qu'un anneau admet toujours un neutre pour la multiplication. Certains ouvrages parlent dans ce cas d'anneau unitaire.} est un triplet \( (A,+,\cdot)\) avec les conditions
    \begin{enumerate}
        \item
            \( (A,+)\) est un groupe abélien. Nous notons \( 0\) le neutre.
        \item
            La multiplication est associative et nous notons \( 1\) le neutre.
        \item
            La multiplication est distributive par rapport à l'addition.
    \end{enumerate}
\end{definition}

\begin{definition}
    Le \defe{centralisateur}{centralisateur} de \( x\in A\) dans \( A\) est l'ensemble
    \begin{equation}
        \{ y\in A\tq xy=yx \},
    \end{equation}
    le \defe{centre}{centre!d'un anneau} de \( A\) est
    \begin{equation}
        \{ y\in A\tq xy=yx,\forall x\in A \}.
    \end{equation}
\end{definition}

\begin{definition}
    Nous disons que \( A\) est un \defe{anneau commutatif}{anneau commutatif} si pour tout \( a,b\in A\) nous avons \( ab=ba\).
\end{definition}

\begin{lemma}       \label{LEMooVUSMooWisQpD}
    Pour tout élément \( a\) d'un anneau nous avons \( a\cdot 0=0\).
\end{lemma}

\begin{proof}
    L'élément \( 0\) est le neutre de l'addition. Il peut être écrit \( 1-1\), et en utilisant la distributivité,
    \begin{equation}
        a\cdot 0=a\cdot (1-1)=a-a=0.
    \end{equation}
    Notons que la dernière égalité s'écrit en détail \( a-a=a+(-a)\) qui donne le neutre de l'addition.
\end{proof}

\begin{proposition}     \label{PROPooNCCGooXjVyVt}
    Dans un anneau non nul, le neutre pour l'addition est distinct du neutre pour la multiplication.
\end{proposition}
\begin{proof}
    Supposons par contraposée que dans un anneau $A$, \( 1 = 0 \). Alors, pour tout \( a \in A \), on a \( a = 1a = 0a = (1 - 1)a = a - a=0 \), d'où l'on déduit \( -a = 0  \) et par suite, \( a = 0. \)
\end{proof}

\begin{normaltext}
Soit \( X\) un ensemble et un anneau $(A, +, \times)$. Nous considérons \( \Fun(X,A)\)\nomenclature[A]{\( \Fun(X,Y)\)}{les applications de \( X\) vers \( Y\)} l'ensemble des applications \( X\to A\). Cet ensemble devient un anneau avec les définitions
\begin{subequations}
    \begin{align}
        (f+g)(x)=f(x)+g(x)\\
        (fg)(x)=f(x)g(x).
    \end{align}
\end{subequations}
Cela est la \defe{structure canonique}{structure d'anneau canonique} d'anneau sur \( \Fun(X,A)\).
\end{normaltext}


\subsection{Morphismes d'anneaux, sous-anneaux, idéaux}
%------------------------------------------------------

\begin{definition}[Morphisme d'anneaux\cite{ooZRUJooXyxPqQ}]      \label{DEFooSPHPooCwjzuz}
    Si \( (A,+,\cdot)\) et \( (B,+,\cdot)\) sont des anneaux, un \defe{morphisme d'anneaux}{morphisme!d'anneaux} est une application \( f\colon A\to B\) telle que
    \begin{enumerate}
        \item \( f(a+b)=f(a)+f(b)\)
        \item
            \( f(a\cdot b)=f(a)\cdot f(b)\)
        \item
            \( f(1)=1\).
    \end{enumerate}
    Étant bien entendu que les significations de \( 1\), $+$ et \( \cdot\) sont différentes à gauche et à droite.
\end{definition}

Remarquons que si \( f\) est un morphisme, nous avons \( f(0)=0\) et \( f(x)^{-1}=f(x^{-1})\).

\begin{definition}
    Un \defe{isomorphisme d'anneaux}{isomorphisme!d'anneaux} est un morphisme bijectif.
\end{definition}

\begin{definition}  \label{DefAJVTPxb}
    Un sous-ensemble \( B\subset A\) d'un anneau est un \defe{sous anneau}{sous anneau} si
    \begin{enumerate}
        \item
            \( 1\in B\)
        \item
            \( B\) est un sous-groupe pour l'addition
        \item
            \( B\) est stable pour la multiplication.
    \end{enumerate}
\end{definition}

\begin{definition}[Idéal dans un anneau]  \label{DefooQULAooREUIU}
    Un sous-ensemble \( I\subset A\) est un \defe{idéal à gauche}{idéal!dans un anneau} à gauche si
    \begin{enumerate}
        \item
            \( I\) est un sous-groupe pour l'addition,
        \item
            pour tout \( a\in A\), \( aI\subset I\).
    \end{enumerate}

    Lorsqu'un ensemble est idéal à gauche et à droite, nous disons que c'est un \defe{idéal bilatère}{idéal!bilatère}. Lorsque nous parlons d'idéal sans précisions, nous parlons d'idéal bilatère.
\end{definition}

\begin{remark}
    Un idéal n'est pas toujours un anneau parce que l'identité pourrait manquer. Un idéal qui contient l'identité est l'anneau complet.
\end{remark}

\begin{definition}[Idéal engendré par un élément]  \label{DefSKTooOTauAR}
    Si \( p\) est un élément d'un anneau \( A\) alors nous notons \( (p)\)\nomenclature[A]{\( (p)\)}{idéal engendré par \( p\)}\index{engendré!idéal dans un anneau} l'idéal dans \( A\) \defe{engendré}{engendré} par \( p\), c'est-à-dire \( pA\).
\end{definition}

\begin{example}
    L'ensemble \( 2\eZ\) est un idéal de \( \eZ\). On peut aussi le noter \( (2) \).
\end{example}

\begin{propositionDef}      \label{PROPooGXMRooTcUGbi}
    Soit \( A\), un anneau, \( I\) un idéal bilatère\footnote{Définition~\ref{DefooQULAooREUIU}.} de \( A\). Nous considérons la relation d'équivalence \( x\sim y\) si et seulement si \( x-y\in I\). Sur le quotient\footnote{Définition \ref{DEFooRHPSooHKBZXl}.}
    \begin{equation}
        A/\sim=A/I,
    \end{equation}
    nous mettons les opérations
    \begin{enumerate}
        \item
            \( [x]+[y]=[x+y]\)
        \item
            \( [x][y]=[xy]\).
    \end{enumerate}
    Nous avons alors les résultats suivants :
    \begin{enumerate}
        \item       \label{ITEMooEJPEooRKAqmS}
            Les opérations sont bien définies;
        \item       \label{ITEMooYBEGooTlHgNz}
            l'ensemble \( A/I\), muni de ces opérations, est un anneau;
        \item       \label{ITEMooLNRLooMkoWXZ}
            la surjection canonique \( \pi\colon A\to A/I\) est un morphisme.
    \end{enumerate}
    Cet anneau est appelé \defe{anneau quotient}{anneau!quotient par un idéal}.
\end{propositionDef}

\begin{proof}
    En plusieurs parties.
    \begin{subproof}
        \item[Pour \ref{ITEMooEJPEooRKAqmS}]
            Nous savons que, par définition,
            \begin{equation}
                \bar x=\{ x+i\tq i\in I \}.
            \end{equation}
            Calculons le produit de représentants génériques de \( \bar x\) et de \( \bar y\) :
            \begin{equation}
                (x+i_1)(y+i_2)=xy+xi_2+yi_1+i_1i_2.
            \end{equation}
            Vu que \( I\) est un idéal nous avons \( xi_2+yi_1+i_1i_2\in I\) et donc bien
            \begin{equation}
                (x+i_1)(y+i_2)\in \overline{ xy }.
            \end{equation}
        \item[Pour \ref{ITEMooYBEGooTlHgNz}]
            Il s'agit de vérifier les conditions de la définition \ref{DefHXJUooKoovob}.

            Vu que \( (A,+)\) est un groupe abélien de neutre \( 0\), nous avons
            \begin{equation}
                [a]+[b]=[a+b]=[b+a]=[b+a],
            \end{equation}
            ainsi que
            \begin{equation}
                [a]+[0]=[a+0]=[a]
            \end{equation}
            et
            \begin{equation}
                [a]+([b]+[c])=[a]+[b+c]=[a+b+c]=[a+b]+[c].
            \end{equation}
            Donc \( (A/I,+)\) est un groupe abélien de neutre \( [0]\).

            L'associativité de \( A\) donne l'associativité dans \( A/I\) :
            \begin{equation}
                \big( [a][b] \big)[c]=[ab][c]=[abc]=[a][bc]=[a]\big( [b][c] \big).
            \end{equation}
            Et enfin pour la distributivité,
            \begin{equation}
                [a]\big( [b]+[c] \big)=[a][b+c]=[a(b+c)]=[ab+ac]=[ab]+[ac]=[a][b]+[a][c].
            \end{equation}
            Nous avons prouvé que \( A/I\) est un anneau de neutre \( [0]\) et d'unité \( [1]\).
        \item[Pour \ref{ITEMooLNRLooMkoWXZ}]
            Nous devons vérifier les trois conditions de la définition \ref{DEFooSPHPooCwjzuz}. Cela est immédiat parce que \( \pi(x)=[x]\).
    \end{subproof}
\end{proof}

\begin{proposition}[Premier théorème d'isomorphisme pour les anneaux]
    Soient \( A\) et \( B\) des anneaux et un homomorphisme \( f\colon A\to B\). Nous considérons l'injection canonique \( j\colon f(A)\to B\) et la surjection canonique \( \phi\colon A\to A/\ker f\). Alors il existe un unique isomorphisme
    \begin{equation}
        \tilde f \colon A/\ker f\to f(A)
    \end{equation}
    tel que \( f=j\circ\tilde f\circ\phi\).

    \begin{equation}
        \xymatrix{%
        A \ar[r]^{f}\ar[d]_{\phi}        &   B\ar[d]^{j}\\
           A/\ker f \ar[r]_{\tilde f}   &   f(A)\subset B
           }
    \end{equation}
\end{proposition}
\index{théorème!isomorphisme!premier!pour les anneaux}

\begin{proposition}     \label{PropIJJIdsousphi}
    Soient \( I\) un idéal de \( A\) et la projection canonique
    \begin{equation}
        \phi\colon A\to A/I.
    \end{equation}
    Elle est une bijection entre les idéaux de \( A\) contenant \( I\) et les idéaux de \( A/I\).

    Dit de façon imagée :
    \begin{equation}        \label{EqKbrizu}
        \{ \text{idéaux de } A\text{ contenant } I\}\simeq\{ \text{idéaux de } A/I \}.
    \end{equation}
\end{proposition}

\begin{proof}
    Si \( I\subset J\) et si \( J \) est un idéal de \( A\), alors \( \phi(J)\) est un idéal dans \( A/I\). En effet un élément de \( \phi(J)\) est de la forme \( \phi(j)\) et un élément de \( A/I\) est de la forme \( \phi(i)\). Leur produit vaut
    \begin{equation}
        \phi(i)\phi(j)=\phi(ij)\in\phi(J).
    \end{equation}

    Soit maintenant \( K\) un idéal dans \( A/I\) et soit \( J=\phi^{-1}(K)\). Étant donné qu'un idéal doit contenir \( 0\) (parce qu'un idéal est un groupe pour l'addition), \( [0]\in K\) et par conséquent \( I\subset\phi^{-1}(K)\).
\end{proof}
% TODO : il faudrait dire à peu près ici qu'une des utilités de Z_2 est le groupe modulaire PSL(2,Z)=SL(2,Z)/Z_2



\subsection{Diviseurs dans un anneau}
%------------------------------------

\begin{definition}[Diviseurs dans un anneau]\label{DiviseursAnneau}
	Soient \( a, b \in A \). On dit que $a$ divise $b$, ou que $a$ est un \defe{diviseur (à gauche)}{diviseur!dans un anneau} de $b$ s'il existe \( c \in A \) tel que \( ac = b \). On dit que c'est un diviseur de $b$ à droite si \( ca = b \) pour un certain \( c \in A \).
	
	Quand $a$ divise $b$, on note \( a\divides b\).
	
	 Si \( S\) est une partie de \( A\), nous notons \( a\divides S\) pour exprimer que \( a\divides x\) pour tout \( x\in S\); de la même façon, \( S\divides b\) signifie que \( x\divides b\) pour tout \( x\in S\).
\end{definition}

\begin{definition}\label{DefrYwbct}
    Soient \( A\) un anneau commutatif et \( S\subset A\). Nous disons que \( \delta\in A\) est un \defe{PGCD}{pgcd!dans un anneau intègre} de \( S\) si
    \begin{enumerate}
        \item
            \( \delta\) divise tous les éléments de \( S\).
        \item
            si \( d\) divise également tous les éléments de \( S\), alors \( d\) divise \( \delta\).
    \end{enumerate}
    Nous disons que \( \mu\in A\) est un \defe{PPCM}{ppcm!dans un anneau intègre} de \( S\) si
    \begin{enumerate}
        \item
            tous les éléments de \( S\) divisent \( \mu\),
        \item
            si tous les éléments de \( S\) divisent un élément \( m \in A\) , alors \( \mu \) divise aussi \( m\).
    \end{enumerate}
\end{definition}

\begin{remark}
    Au sens de la définition \ref{DefrYwbct}, le pgcd n'est pas unique. Dans \( \eZ\) par exemple les nombres \( 4\) et \( -4\) sont tout deux pgcd de \( \{4,16  \}\).

    Dans \( \eZ\) cependant, nous modifions implicitement la définition et nous n'acceptons que les positifs, de telle sorte à ce que l'unique pgcd soit effectivement le plus grand pour l'ordre usuel sur \( \eZ\).

    Pour l'unicité dans \( \eZ\), voir \ref{LEMooBJVJooFyuFeN}.
\end{remark}


\begin{definition}
Un élément \( a\in A\) est \defe{régulier à droite}{régulier à droite} si \( ba=0\) implique \( b=0\). Il est régulier à gauche si \( ab=0\) implique \( b=0\).
\end{definition}

\subsection{Nilpotents, unipotents, inversibles}
%----------------------------------------------

\begin{definition}[Éléments nilpotents, unipotents et inversibles]
	On dit que \( a \in A \) est \defe{nilpotent}{nilpotent} s'il existe \( n \in \eN \) tel que \( a^n = 0 \). Il est dit \defe{unipotent}{unipotent} si \( a-1\) est nilpotent, c'est-à-dire si \( (a-1)^n =0\) pour un certain \( n \in \eN \).

	Un élément \( a \in A \) est dit \defe{inversible}{élément!inversible!dans un anneau} s'il existe \( b \in A \) tel que \( ab = 1 \).
\end{definition}

L'ensemble \( U(A)\)\nomenclature[A]{\( U(A)\)}{ensemble des inversibles} des éléments inversibles de \( A\) est un groupe pour la multiplication. Nous notons \( A^*=A\setminus\{ 0 \}\).

\begin{lemma}
    Si \( a\) et \( b\) commutent, nous avons, pour tout \( r \in \eN \) et \( r > 0\), la formule
    \begin{equation}        \label{Eqarpurmkbk}
        a^{r+1}-b^{r+1}=(a-b)\left(\sum_{k=0}^ra^{r-k}b^k\right).
    \end{equation}
\end{lemma}

\begin{proof}
  Démontrons cela par récurrence. Le cas \( r = 0 \) est évident. Pour
  un \( r \) donné, si \eqref{Eqarpurmkbk} est vraie, alors
  \begin{align*}
    a^{r+2}-b^{r+2}&= a^{r+1}a - a^{r+1}b +a^{r+1}b - b^{r+1}b\\
    &= a^{r+1}(a - b) + (a^{r+1} - b^{r+1})b\\
    &= a^{r+1}(a - b) + (a-b)\left(\sum_{k=0}^ra^{r-k}b^k\right)b\\
    &= (a - b) \left(a^{r+1} + \left(\sum_{k=0}^ra^{r-k}b^k\right)b\right)\\
    &= (a - b) \left(a^{r+1} + \sum_{k=0}^ra^{r-k}b^{k + 1}\right)\\
    &= (a - b) \left(a^{r+1} + \sum_{k'=1}^{r+1}a^{(r+1)-k'}b^{k'}\right)\\
    &= (a - b) \left(\sum_{k'=0}^{r+1}a^{(r+1)-k'}b^{k'}\right).
  \end{align*}
\end{proof}

\begin{proposition}
    Si \( a\) est un élément nilpotent de l'anneau \( A\), alors \( 1-a\) est inversible. Si \( a\) est nilpotent non nul, alors il est diviseur de zéro.
\end{proposition}

\begin{proof}
    Soit \( n\) le minimum tel que \( a^n=0\). En vertu de la formule \eqref{Eqarpurmkbk} nous avons
    \begin{equation}
        1=1-a^n=(1-a)(1+a+\cdots+a^{n-1})=(1+a+\cdots+a^{n-1})(1-a).
    \end{equation}
    La somme \( 1+a+\cdots+a^{n-1}\) est donc un inverse de \( (1-a)\).
\end{proof}


\subsection{Diviseurs de zéro et intégrité}
%--------------------------------------------

Conformément à la définition \ref{DiviseursAnneau} de diviseur, nous posons la définition suivante pour les diviseurs de zéro.
\begin{definition}[\cite{ooTNKJooSCSCZQ}]
    Un élément \( a\neq 0\) est un \defe{diviseur de zéro à gauche}{diviseur!de zéro} s'il existe \( x\neq 0\) tel que $xa=0$. L'élément \( a\) est un diviseur de zéro \defe{à droite}{diviseur!de zéro à droite} s'il existe \( b\) tel que \( ab=0\).

    Nous disons que \( a\) est un \defe{diviseur de zéro}{diviseur de zéro} si il est un diviseur de zéro à gauche ou à droite.
\end{definition}

L'absence de tels diviseurs dans un anneau est une propriété intéressante: on dit dans ce cas que l'anneau est intègre. Nous étudions ces anneaux plus en détail en section~\ref{SECAnneauxIntegres}. D'ores et déjà, soulignons des propriétés équivalentes.

\begin{propositionDef}[\cite{MonCerveau}]           \label{DEFooTAOPooWDPYmd}
    Soit $A$ un anneau non réduit à \( \{ 0 \}\). Les assertions suivantes sont équivalentes:
    \begin{enumerate}
        \item       \label{ITEMooMXMKooXMYpkN}
            \( A\) ne possède pas de diviseurs de zéro.
        \item       \label{ITEMooLAJCooFwxXrV}
            La règle du produit nul s'applique dans $A$: pour tous \( a, b \in A \), si \( ab=0\), alors \( a = 0\) ou \( b = 0\).
            \index{règle du produit nul}
        \item       \label{ITEMooQNTFooSRrVPK}
            On peut simplifier par un même élément non-nul, deux expressions produit dans $A$ qui sont égales: pour tous \( a, b, c \in A \) avec \( a \neq 0 \), si \( ab = ac \), alors \( b = c \).
    \end{enumerate}
    Un anneau non réduit à \( \{ 0 \}\) qui vérifie ces propriétés est dit \defe{intègre}{anneau intègre}.
\end{propositionDef}

\begin{proof}
    En trois implications.
    \begin{subproof}
        \item[\ref{ITEMooMXMKooXMYpkN} implique \ref{ITEMooLAJCooFwxXrV}]

            Si \( ab=0\) avec \( a,b\neq 0\) alors \( a\) est un diviseur de zéro. Vu que nous supposons que \( A\) n'a pas de diviseurs de zéro, soit \( a\) soit \( b\) est nul.
        \item[\ref{ITEMooLAJCooFwxXrV} implique \ref{ITEMooQNTFooSRrVPK}]

            Si \( ab=ac\), alors \( a(b-c)=0\) et l'hypothèse dit que soit \( a=0\), soit \( b-c=0\). Donc si \( a\neq 0\), alors \( b-c=0\), et donc \( b = c \).
        \item[\ref{ITEMooQNTFooSRrVPK} implique \ref{ITEMooMXMKooXMYpkN}]
            Si \( A=\{ 0 \}\), le point \ref{ITEMooQNTFooSRrVPK} n'est pas applicable.

            Si \( a\neq 0\) et \( xa=0\), alors nous avons aussi \( xa=0\times a\). Par propriété de simplification, \( x=0\). Donc \( a\) n'est pas un diviseur de zéro à gauche. Nous prouvons de la même façon qu'il n'y a pas de diviseurs de zéro à droite.
    \end{subproof}
\end{proof}

Voici une autre vision de la divisibilité dans les anneaux intègres.
\begin{proposition}\label{PropDiviseurIdeaux}
    Dans un anneau intègre\footnote{Définition \ref{DEFooTAOPooWDPYmd}.} $A$, on a l'équivalence suivante concernant deux éléments \( a, b \in A \):
\begin{equation}
    a\divides b\Leftrightarrow (b)\subset (a).
\end{equation}
\end{proposition}

Dans les anneaux intègres; la divisibilité devient en réalité une relation d'ordre dont nous pouvons chercher un maximum et un minimum.

\begin{lemma}[\cite{ooLIOMooBuCPUS}]
    Les éléments inversibles d'un anneau sont diviseurs de tous les éléments.
\end{lemma}

\begin{proof}
    Soit \( k\) inversible d'inverse \( k'\) : \( kk'=1\); soit aussi \( a\in A\). Alors \( a=k(k'a)\), ce qui montre que \( k\) divise \( a\).
\end{proof}

\begin{lemma}[\cite{ooLIOMooBuCPUS}]
    Dans un anneau, \( 1\) est un pgcd de \( a\) et \( b\) si et seulement si les seuls diviseurs communs sont les inversibles.
\end{lemma}

\begin{proof}
    Supposons pour commencer que \( 1\) est un pgcd de \( a\) et \( b\). Un diviseur commun de \( a\) et \( b\) doit donc diviser \( 1\). Or un diviseur de \( 1\) est forcément inversible.

    Dans l'autre sens, les diviseurs communs de \( a\) et \( b\) sont tous inversibles et donc diviseurs de \( 1\). Donc \( 1\) est un pgcd de \( a\) et \( b\).
\end{proof}


%+++++++++++++++++++++++++++++++++++++++++++++++++++++++++++++++++++++++++++++++++++++++++++++++++++++++++++++++++++++++++++ 
\section{Somme à valeurs dans un groupe abélien}
%+++++++++++++++++++++++++++++++++++++++++++++++++++++++++++++++++++++++++++++++++++++++++++++++++++++++++++++++++++++++++++

%--------------------------------------------------------------------------------------------------------------------------- 
\subsection{Somme à valeurs dans un anneau}
%---------------------------------------------------------------------------------------------------------------------------

\begin{definition}      \label{DEFooNEVNooJlmJOC}
    Si \( f\colon \{ 0,\ldots, N \}\to A\) est une application vers un anneau \( A\), alors nous définissons la notation \( \sum_{i=0}^Nf(i)\) par récurrence de la façon suivante :
    \begin{enumerate}
        \item
            \( \sum_{i=0}^0f(i)=f(0)\),
        \item
            \( \sum_{i=0}^{k}f(i)=\sum_{i=0}^{k-1}f(i)+f(k)\).
    \end{enumerate}
\end{definition}

%--------------------------------------------------------------------------------------------------------------------------- 
\subsection{Somme à valeurs dans un groupe abélien}
%---------------------------------------------------------------------------------------------------------------------------

Si \( S\) est un ensemble fini, nous savons de la proposition \ref{PROPooJLGKooDCcnWi} qu'il existe un unique \( N\in \eN\) pour lequel il existe une bijection \( \varphi\colon \{ 0,\ldots, N \}\to S\). Cette bijection n'est à priori pas unique.

\begin{lemmaDef}[\cite{MonCerveau}]       \label{DEFooLNEXooYMQjRo}
    Soient un groupe abélien \( (G,+)\) ainsi qu'un ensemble fini \( I\) contenant \( n\) éléments. Soit une application \( f\colon I\to G \). Si \( \sigma_1,\sigma_2\colon \{1,\ldots, n \}\to I\) sont deux bijections, alors\footnote{Pour rappel, le symbole \( \sum_{i=1}^n\) est défini par \ref{DEFooNEVNooJlmJOC}.}
    \begin{equation}
        \sum_{i=1}^nf\big( \sigma_1(i) \big)=\sum_{i=1}^nf\big( \sigma_2(i) \big).
    \end{equation}
    La valeur commune est notée
    \begin{equation}
        \sum_{i\in I}f(i)
    \end{equation}
\end{lemmaDef}

La définition \ref{DEFooLNEXooYMQjRo} donne lieu à un certain nombre de remarques.
\begin{enumerate}
    \item
        Elle donne la somme sur un ensemble fini. Un problème avec les ensembles infinis (outre la convergence) est l'ordre de sommation. Si vous voulez sommer sur \( \eZ\), dans quel ordre le faire ?
    \item
        Pour aller plus loin, et sommer sur des ensembles infinis, il faut regarder la définition \ref{DefHYgkkA}. 
\end{enumerate}

\begin{proposition}     \label{PROPooJBQVooNqWErk}
    Soient un groupe abélien \( (G,+)\), un ensemble fini \( I\) est un ensemble fini, une application \( f\colon I\to G\) et une bijection \( \sigma\colon I\to I\). Alors
    \begin{equation}
        \sum_{i\in I}f(i)=\sum_{i\in I}f\big( \sigma(i) \big).
    \end{equation}
\end{proposition}

Si nous avons une application \( L\colon S\to S\), nous notons
\begin{equation}
    \sum_{s\in S}f\big( L(s) \big)=\sum_{s\in S}(f\circ L)(s).
\end{equation}
Cette façon d'écrire donne une interprétation pour la notation \( \sum_{g\in G}f(hg)\) qui arrive dans la proposition \ref{PROPooWJQQooFINSEc}. Il s'agit de considérer l'application \( L_h\) du lemme \ref{LEMooBIBFooBHxFYC}, de considérer\footnote{Le fait que \( L_h\) soit une bijection n'a pas d'importance ici.}
\begin{equation}        \label{EQooQQBEooFDOBVG}
    \sum_{g\in G}f(hg)=\sum_{g\in G}(f\circ L_h)(g)
\end{equation}
et de faire tourner la définition \ref{DEFooLNEXooYMQjRo}. La même chose tient pour définir \( \sum_{g\in G}(gh)\) à l'aide de \( R_h\).

\begin{proposition}[\cite{MonCerveau}]     \label{PROPooWJQQooFINSEc}
    Soient un groupe fini \( G\) et une fonction \( f\colon G\to A\) où \( A\) est un anneau commutatif. Alors
    \begin{equation}
        \sum_{g\in G}f(g)=\sum_{g\in G}f(gh)=\sum_{g\in G}f(hg)
    \end{equation}
    pour tout \( h\in G\).
\end{proposition}

\begin{proof}
    Nous avons une bijection \( \varphi\colon \{ 0,\ldots,  N \}\to G\) garantie par la proposition \ref{PROPooJLGKooDCcnWi}. La définition est que
    \begin{equation}
        \sum_{g\in G}f(g)=\sum_{i=0}^Nf\big( \varphi(i) \big).
    \end{equation}
    Par ailleurs, le lemme \ref{LEMooBIBFooBHxFYC} donne une bijection \( L_h\colon G\to G\) et permet de considérer la composée
    \begin{equation}
        \begin{aligned}
            \varphi'\colon \{ 0,\ldots,  N \}&\to G \\
            \varphi'=L_h\circ \varphi.
        \end{aligned}
    \end{equation}
    La proposition \ref{DEFooLNEXooYMQjRo} nous permet d'utiliser la bijection \( \varphi'\) au lieu de \( \varphi\) pour exprimer la somme \( \sum_{g\in G}\). Ensuite un jeu de notation utilisant \eqref{EQooQQBEooFDOBVG} donne
    \begin{equation}
        \begin{aligned}[]
            &\sum_{g\in G}f(g)=\sum_{i=0}^Nf\big( \varphi(i) \big)=\sum_{i=0}^Nf\big( \varphi'(i) \big)=\sum_{i=0}^N(f\circ L_h\circ \varphi)(i)\\
            &\quad=\sum_{i=0}^N(f\circ L_h)\big( \varphi(i) \big)=\sum_{g\in G}(f\circ L_h)(g)=\sum_{g\in G}f(hg).
        \end{aligned}
    \end{equation}
    En ce qui concerne \( \sum_{g\in G}f(gh)\), c'est la même chose, en utilisant \( R_h\) au lieu de \( L_h\).
\end{proof}

Tout cela nous permet de faire une somme sympathique et bien connue.
\begin{lemma}
    Soit \( n\in \eN\). Nous avons
    \begin{equation}
        \sum_{k=0}^nk=\frac{ n(n+1) }{ 2 }.
    \end{equation}
\end{lemma}

\begin{proof}
    La preuve est pratiquement immédiate par récurrence. Nous allons donner une preuve plus «constructive», qui formalise l'idée classique d'écrire la somme à l'endroit et à l'envers.


    Nous notons \( S\) la somme \( \sum_{k=0}^nk\). Le lemme \ref{DEFooLNEXooYMQjRo} dit que si \( \sigma_i\colon \{ 0,\ldots, n \}\to \{ 0,\ldots, n \}\) sont deux bijections, alors \( \sum_{k=0}^nf\big( \sigma_1(k) \big)=\sum_{k=0}^nf\big( \sigma_2(k) \big)\). Nous sommes intéressé au cas \( f(i)=i\).

    En prenant \( \sigma_1(k)=k\) et \( \sigma_2(k)=n-k\), nous avons
    \begin{equation}
        S=\sum_{k=0}^nk=\sum_{k=0}^n(n-k).
    \end{equation}
    Donc
    \begin{equation}
        2S=\sum_{k=0}^n\big( k+(n-k) \big)=\sum_{k=0}^nn=n\sum_{k=0}^n1=n(n+1).
    \end{equation}
    En divisant par deux, nous obtenons le résultat annoncé.
\end{proof}

Dans le même ordre d'idée, pour le produit.
\begin{proposition}     \label{PROPooQMUDooQQVRIe}
    Si \( E\) est un ensemble fini et si \( G\) est un groupe abélien, alors pour toute fonction \( f\colon E\to G\) et pour toute permutation \( \sigma\) de \( E\),
    \begin{equation}
        \prod_{i\in E}f(i)=\prod_{i\in E}f\big( \sigma(i) \big)
    \end{equation}
\end{proposition}


